\section{Testing Evidence}\label{sec:test-evidence}

Quantitative testing claims in this chapter rest on the metrics baseline collected on 2026-06-01 at monorepo commit
\texttt{7e8a55fb}.
Section~\ref{sec:testing-ci} documents the full CI matrix; here we summarize the numbers most relevant to evaluation
claims.

\begin{table}[htbp]
  \centering
  \footnotesize
  \renewcommand{\arraystretch}{1.25}
  \begin{tabularx}{\textwidth}{>{\raggedright\arraybackslash}p{2.4cm}rr>{\raggedright\arraybackslash}X}
    \toprule
    \textbf{Subsystem} & \textbf{Tests} & \textbf{Pass (local)} & \textbf{Coverage (where verified)} \\
    \midrule
    Compiler (CTest) & 635 registered & 635 / 635 & gcovr in CI; local \texttt{lcov} unavailable \\
    Server (Vitest) & 31 cases & 20 unit (no Postgres); 11 integration skipped locally & 34.4\% lines (unit-only) \\
    Client (Vitest) & 79 cases & 79 / 79 & 51.4\% lines \\
    \bottomrule
  \end{tabularx}
  \caption{Verified test inventory for evaluation claims (2026-06-01).}\label{tab:evaluation-test-counts}
\end{table}

The compiler suite spans parsing, explicit type rules, lowering, MIDI writing, integration through
\texttt{motivo::compile()}, and four golden \texttt{.motivo} fixtures that assert stage-specific diagnostics or MIDI
shape.
During baseline collection, the golden test \texttt{Golden.LoweringErrorsAreCollectedAcrossTracks} initially failed
because \texttt{lowering\_errors.motivo} relied on integer division semantics that changed when \texttt{/} began
promoting to \texttt{double}.
The fixture was corrected to trigger lowering errors via invalid loop bounds; the full suite now passes at \textbf{635 /
635}.

Coverage percentages are not a proof of correctness.
They do not show whether the language is pleasant to use or whether every musical edge case has been modeled.
They do show that the implementation is treated as a software system with repeatable checks at compiler, API,
client, and
paper boundaries.
For this thesis, that matters because Motivo is evaluated both as a language design and as a full-stack engineering
project.
